% Title: RAIMD:

\documentclass[11pt]{article}
\usepackage{latex8}
\usepackage{times}
\usepackage{graphicx}
\usepackage{subcaption}
\usepackage{url}
\usepackage{amsmath,amsfonts}
\usepackage{xspace}
\usepackage{setspace}


% NSF required minimum margin is 1 inch
\usepackage[margin=1.0in,includefoot]{geometry}
%\linespread{1.15}

\newcommand{\DDRx}{DDR{\em x}\xspace}
%\newcommand{\mCore}{{\em m}Core\xspace} \newcommand{\tAA}{$^\text{t}$AA\ }
%\newcommand{\tWAL}{$^\text{t}$WAL\ }
\newcommand{\tAA}{tAA\xspace} \newcommand{\tWAL}{tWAL\xspace}
\newcommand{\tRFC}{tRFC\xspace}

\newcommand{\eg}{e.g.\xspace} \newcommand{\ie}{i.e.\xspace}
\newcommand{\etal}{et al.\xspace} \newcommand{\naive}{na\"ive\xspace}

\newcommand{\Fig}[1]{Figure~\ref{fig:#1}\xspace}
\newcommand{\Sec}[1]{Section~\ref{sec:#1}\xspace}
\newcommand{\Alg}[1]{Algorithm~\ref{alg:#1}\xspace}
\newcommand{\Tab}[1]{Table~\ref{tab:#1}\xspace}
\newcommand{\Eq}[1]{Equation~\eqref{eq:#1}\xspace}
%\newcommand{\Para}[1]{\vspace{0.04in}{\bf #1}\xspace}
\newcommand{\Para}[1]{\paragraph{#1}}

\newcommand{\comment}[1]{} \newcommand{\note}[1]{[{\em #1}]}
\newcommand{\ZZ}[1]{[{\em #1}]} \newcommand{\todo}[1]{} \newcommand{\zzc}[1]{}
\newcommand{\zz}[1]{} \newcommand{\zzcite}[1]{}

% add "C" to section heading \renewcommand{\thesection} {C.\arabic{section}}
% \renewcommand{\thepage} {C-\arabic{page}}

\makeatletter \newcommand\figcaption{\def\@captype{figure}\caption}
\newcommand\tabcaption{\def\@captype{table}\caption} \makeatother

\begin{document}

% Proposal Title:
%\centerline{\large\bf RAIMD: Highly Reliable Memory Architecture for Next
%Generation Device Technologies}
%\centerline{\large\bf RAIMD: Versatile and Highly Reliable Memory Architecture
%for Next Generation Memories}

\begin{center}
{\Large\bf SHF: Energy-Efficient and Cost-Effective Architecture for
Highly Reliable Memory Systems with Current and Emerging Technologies}
\end{center}
\vspace{0.5cm}

% Here me want to discuss the motivation based on challenges in future memory

\Section{Research Overview}

% DRAM errors call for stronger protection "was getting worse" -- it is too
% casual Plan: read the papers again and sharpen the writing

Memory reliability has become a severe concern for contemporary memory systems.
Two large-scale field studies in 2009~\cite{2009-sigm-dram-error} and
2012~\cite{2012-sc-dram-error} reported that DRAM memories had much
higher raw error rates than expected.  More importantly, DRAM errors were found
to be highly correlated. The findings raised a serious question about the
effectiveness of memory error protection schemes used in product systems,
because those schemes are designed against random errors.  A later study in
2015~\cite{2015-asplos-dram-error} not only confirmed the previous findings but
also showed a worrisome trend of increasingly weakening DRAM reliability because
of the advance in DRAM fabrication technology.  The technology advance has
drastically improved DRAM density, bandwidth, and energy efficiency, but
unfortunately it has weakened DRAM cell reliability because of shrunk cell size.

% DRAM errors causes security issues

The weakening memory reliability also has a serious consequence: it incurs a new type
of vulnerability in computer security.  The standing out issue is Rowhammer
attack~\cite{kim2014flipping,mutlu2017rowhammer}, which exploits the weakening cell reliability of high-density
DRAM devices.  In such an attack, an adversary runs a malicious program to flip
data bits in memory; and with a carefully crafted program, the adversary can
obtain privileged access to the system~\cite{Seaborn}.  This is a severe problem and
surprisingly challenging to address.  Many designs have been proposed, and JEDEC
has made an extension to the DDR5 standard~\cite{JEDEC:2024}, to address the problem.
However, it persists and even worsens, because DRAM cell reliability continues
to weaken as DRAM fabrication technology advances~\cite{Canpolat:hpca25}.  The latter is a
bless for memory capacity, bandwidth, and energy efficiency, but not for
reliability.

% Check "standing out" with my new ipad．It's very good.

% What we propose. Stay at high-level to make the picture clear. Save challenges
% to the next paragraph.

The PI and co-PI propose a new memory architecture, called RAIMD (Redundant
Array of Independent Memory Device), and will fully investigate the new
architecture to address those challenges.  RAIMD is similar to RAID (Redundant
Array of Inexpensive Disks) in organization: memory devices in RAIMD are
independent with local access scheduling, contain local error correcting codes,
and may form parity groups for error correction.  It supports strong reliability
with low cost and very high energy efficiency.  It is also highly flexible: a
RAIMD memory can be dynamically configured as RAIMD-0, RAIMD-5, RAIMD-6 or other
variants as done in RAID, with varying ratio of redundancy, to allow trade-offs
between reliability and cost.  It can offer much higher reliability than ECC
memory and Chipkill Correct memory for a given storage budget, particularly when
used with memory devices of weak reliability.

% TODO Must check the term "memory interface protocol" Emerging memory may
% require more support for reliability

RAIMD is designed for reliability but can also enhance security and support new
memory technologies.  RAIMD supports strong end-to-end error detection that
covers memory cells, data bus, address bus, and I/O pins and circuits, which
prevents Rowhammer attacks from succeeding.  It uses decoupled memory access
scheduling, which allows memory devices to work semi-independently with the
memory controller, and thus can easily support devices using NVM (non-volatile
memory)
technologies~\cite{Lee:ISCA09,Zhou:ISCA09,Qureshi:ISCA09,Qureshi:HPCA10}.  Those
technologies have been studied intensively for two decades; and some has been
tried by the industry in product systems but
unsuccessfully~\cite{Waddington:memsys19}.  Because of this property, in the
RAIMD architecture those memory devices can have complex internal circuits to enhance
their durability and reliability, which are critical to their adoption.

% Now it's time to discuss challenges.  Need to mention memory disturbance
% errors in previous paragraph

RAIMD is not a trivial migration of RAID to memory.  There is a critical
difference between memory systems and storage systems: memory latency is
measured in nanoseconds, while disk latency is measured in milliseconds or at
least a fraction of milliseconds.  Any increase of memory latency incurred by
RAIMD could have significant impact on program performance.  We have conducted a
preliminary study to address the performance concern as well as the energy
saving potential of RAIMD.  The initial result has been published in a short
paper~\cite{Wu:cal23}, and a full paper has been submitted and is under review.
The evaluation shows that the average performance overhead of RAIMD-5, which
resembles RAID-5 in organization, is only 3.8\% over a DDR5 baseline for SPEC
CPU2017 and CPU2006 workloads.  It reduces the average memory energy consumption
by 25.9\% from the DDR5 baseline, which has no ChipKill protection.  When
compared to a ChipKill scheme for DDR5, it reduces the average energy
consumption by 69.4\%.  The design is well crafted to eliminate any increase of
memory read latency and to minimize the increase of memory traffic.  It contains
a set of holistic changes to memory controller, device interface, memory command
protocol, memory address mapping, and last-level cache.

RAIMD has several other advantages as a general-purpose, energy-efficient, and
highly reliable memory architecture.  It allows a low-cost implementation
suitable for personal computers and mobile devices: a RAIMD memory can be
constructed using a small number of memory devices, while a conventional
ChipKill design requires many memory devices to be activated to serve one
request.  Its high energy efficiency is not only desired on those computers but
also on servers.  It is also highly configurable: a RAIMD memory can be
configured as RAIMD-0 (no error protection), RAIMD-5, RAIMD-6 (double ChipKill
protection), or other variants; and the configuration is done at booting time,
not manufacturing time.  Therefore, it fits computer systems with very low to
very high reliability requirements, and can be re-configured to just meet
the reliability target.

% TODO Discuss what is decoupled memory access scheduling FIXME We need to
% discuss more about the proposed research in a smooth way

We will conduct a comprehensive research to investigate the full potential of
RAIMD.  In particular, we will study RAIMD with decoupled memory access
scheduling ({\em decoupled scheduling} in short), which has not yet been covered
by the preliminary study.  We believe it is the key to fully addressing the
growing concern on memory reliability and durability, as well as to explore the
opportunities of adopting new memory technologies.  With decoupled scheduling,
the memory controller issues activate, read/write, and other miscellaneous
commands to memory devices, but memory devices can schedule the actual timing of
the corresponding memory operations.   By comparison, with rigid scheduling used
in current memory systems, memory devices must respond to the memory controller
as slave devices.  The change is small but important: A memory device can now be
contain customized enhancement circuits to address any reliability weaknesses
caused by a particular memory technology or fabrication process.  The problems
of DRAM disturbance errors and PCM (Phase-Change-Memory) write endurance can be
fully addressed in this approach.  High-density devices with weakened
reliability can be further enhanced with strong ECC code embedded in the
devices.  When combined with the RAIMD organization, a designer can build a
highly reliable memory system with high energy efficiency and low cost.

A key issue, and the major challenge, is the scheduling of memory bus usage
because the memory controller no longer has the full control of memory devices.
We have conducted another preliminary study on decoupled scheduling for a system
of hybrid DRAM and NVM memory modules~\cite{Trung:pomi}.  The result is
promising: it shows that decoupled scheduling may approach, and even exceed, the
performance of rigid, centralized scheduling.  However, decoupling in RAIMD is
more challenging than in a simple hybrid memory system because the timings of
memory operations in RAIMD will have much higher variance.  We propose to use
perceptron-based neural model to assist decoupled scheduling.  Neural model
based prediction has been shown to be effective for improving the accuracy of
branch prediction~\cite{Jimenez:hpca01,Bhatia:isca19} and memory
prefetching~\cite{Garza:isca19}.  In RAIMD, it will be used to predict memory
operation completion time (and device ready time) with decoupled scheduling,
which we believe is a sound and promising approach.

% double check term "PNN" and perceptron-based neural network

%%% This paragraph is too overwhelming -- it is good, but overwhelming.
%
%Decoupled scheduling is a major deviation from centralized scheduling in a
%conventional memory architecture, and the concern that it may not be as
%efficient as centralized scheduling in memory bandwidth utilization.  We have
%conducted a preliminary study on a different form of decoupled
%scheduling~\cite{} for a hybrid DRAM and NVM memory system.  The result is
%promising: it shows that in a hybrid DDR{\em x} DRAM and NVM memory, decoupled
%scheduling performs very close to centralized scheduling.  The design of
%decoupled scheduling for RAIMD, however, is more challenging than that because
%the timing variance in a RAIMD system is expected to be much higher than that
%in a simple hybrid system.  We propose to use perceptron neural network (NN) to
%improve the efficiency of decoupled scheduling.  Perceptron NN has been shown
%to be effective in branch prediction and memory prefetching.  We believe it is
%promising in RAIMD based on our experience in the preliminary study.

The research agenda is as follows:

\begin{itemize}

\item  To study the integration of decoupled memory access scheduling in the
RAIMD architecture, and the use of device-specific enhancement circuits to
improve reliability, security, and durability of memory devices.

\item To study neural model based prediction to improve the performance of
decoupled scheduling for memory devices of timing diversity.

\item To investigate holistic management of memory reliability, including the
optimal configuration of global protection and local protection for a given
storage budget.

\end{itemize}

\section{Proposed Research}

In this section, we first discuss the characteristics of RAIMD, present two
preliminary studies on a baseline RAIMD design and decoupled memory access
scheduling, respectively, and then discuss the research agenda items in detail.

% Plan: Use the diagram and discussions from CAL paper and TPDS submission
\subsection{Characteristics of RAIMD}

RAIMD has a distinguishing characteristics, when compared to conventional memory
architecture such as DDR4 and DDR5, that a memory request is served by a single
memory device rather than a rank of devices.  A memory block of LLC block size,
typically 64 bytes or 512 bits, is mapped to a single device.  It drastically
reduces memory operational energy, i.e., energy spent on precharge and
activation, at the cost of increased data transfer time.  For example, in a
DDR5-6400 system using x8 devices (with 8-bit data I/O), a memory block of 64
bytes is mapped to four memory devices, and it takes 2.5 ns to transfer the
block on memory bus.  In an equivalent RAIMD system, the block is mapped to a
single device, and thus it takes 10 ns to transfer the block on memory bus.
However, the latency of an idle memory is typically dozens of nanoseconds, and
for memory-intensive workloads the dominant portion of memory latency is
queueing delay.  Thus, the increase of data transfer time may only have a
moderate impact on performance.  In RAIMD, a memory device is effectively a
memory module by itself, and thus memory-level parallelism increases
substantially, which may offset the performance loss due to increased memory
latency.

Because of this characteristics, a RAIMD memory system is flexible in its
organization because it can be configured dynamically during system booting
time.  Therefore, such a system allows trade-offs among performance, energy
efficiency, cost, and reliability.  To one extreme, it can be configured as
RAIMD-0 (no redundancy) if reliability is not a concern for a given workload.
To another extreme, RAIMD-1 (mirroring) offers very high reliability at the cost
of reduced memory capacity.  In practice, RAIMD-5 (with single parity block per
group) and RAIMD-6 (with double parity blocks per group) are of great interest
because they provide high reliability with relatively low cost.  The
parity-to-data ratio of them can be adjusted, giving another option of making
trade-off.  After all, all those RAIMD variants are highly energy-efficient when
compared to a conventional memory architecture of wide rank.

A memory device in RAIMD is required to have extra bits per memory block to
store an error detecting code (EDC).  The requirement is consistent with the
JEDEC DDR5 standard, which requires extra on-die ECC bits in DDR5 devices for
local error detection and correction.  However, EDC in RAIMD is generated and
checked by the memory controller.  Therefore, it is capable of end-to-end error
detection that protects memory cells, I/O pins and circuits, data bus, and
address bus.  Note that for address bus error detection, memory address bits
will be combined with data bits to generate the EDC.
The memory block size in RAIMD is the same as the LLC block size,
which is 512 bits (64 bytes) for most contemporary processors.  Because of the
relatively large block size, EDC in RAIMD is highly efficient.  For example, a
32-bit EDC can provide strong error detection for a 512-bit block, with a
moderate storage overhead of 6.25\%.  A memory device may also contain
additional ECC bits for local error correction (as DDR5 does), which further
enhances the reliability.  We will discus more in
Section~\ref{sec:holistic}.

\subsection{Preliminary Studies}

We have conducted two preliminary studies related to the proposed RAIMD
architecture.  The first study is about a baseline design of RAIMD and its
performance and energy efficiency.  The result is promising: it shows
that with careful design optimizations, RAIMD can support ChipKill protection
with much higher energy efficiency than conventional memory architecture with a
very moderate performance loss (less than 4\%).  The second study is related to
decoupled scheduling of memory access, which is also promising: it shows
decoupled scheduling for DRAM memories can achieve the same or even slightly
better performance than conventional, rigid memory scheduling.

\subsubsection{Baseline RAIMD Design}

\paragraph{Overview.} We have conducted a baseline design of RAIMD based on
DDR5-6400 memory devices.  The RAIMD design contains two critical features to
maintain high memory speed and to reduce memory traffic.  A critical design
issue in RAIMD is the memory layout and address mapping: any complex address
mapping from the physical address to the device address (with channel, rank,
bank, row, and column components) may significantly increase memory latency.
Note that it is not an issue to a storage system, as the latency of a storage device is
orders of magnitude longer than that of a memory device.  Furthermore, a system
designer may prefer a particular address mapping scheme for a sound reason; for
example, to use page-interleaving for spatial locality, cacheline-interleaving for
parallelism, or a variant of their combination for a balance.  This requirement further
complicates the address mapping in RAIMD. Second, since RAIMD uses parity blocks
for error correction, the increase of memory traffic for parity block access is
a concern.  RAID-5 has the small write problem~\cite{RAID}: a write may incur
four disk accesses, namely data block read, parity block read, data block write,
and parity block write.  Fortunately, we have found a solution that not only
eliminates parity block reads (except during error correction) but also
significantly reduces parity writes, which we will discuss in detail.

\paragraph{Address mapping and memory layout.}

% Can we claim this is our scheme?

\begin{figure}[!t] \centering \begin{subfigure}{0.4\textwidth}
\includegraphics[height=4cm]{./graphs/mapping.pdf} \subcaption{Address mapping
and memory layout in the physical address space.} \label{fig:mapping}
\end{subfigure} \hspace{0.5cm} \begin{subfigure}{0.55\textwidth}
\includegraphics[height=4.5cm]{./graphs/organization_no_mapping.png}
\subcaption{Memory layout in the device address space.} \label{fig:raimd_layout}
\end{subfigure} \caption{The framework of address mapping and a generic memory
layout.}
%TODO In the left figure, top: change Data Block to Memory Block.
\label{fig:raimd_overview} \end{figure}

RAIMD solves the problem by using dual address mappings for data blocks and parity blocks, as
shown in Fig.~\ref{fig:raimd_overview}.  Data blocks and parity blocks are
mapped to separate regions of the physical address space with different prefix
address bits.  The address mapping for data blocks only involves bit selection
logic and optionally simple logic operations such as XORing.  Thus, this is no
extra latency for data block access compared to a conventional memory system.
The address mapping for parity blocks, on the other hand, is complex to achieve
a desired block layout.  A Parity Address Mapping logic is applied to the
physical memory address before it is mapped into the device address space
(Fig.~\ref{fig:mapping}).  This logic may add a delay of multiple
nanoseconds to parity block access, but it is not on the critical path of
program execution and thus has negligible performance impact.

\begin{figure*}[!t] \centering
\subfloat[\label{fig:raimd5_address_mapping_stripe}]{%
\includegraphics[width=0.5\linewidth]{./graphs/address_mapping_a.pdf}}
\subfloat[\label{fig:raimd5_mapping_conventional}]{%
\includegraphics[width=0.5\linewidth]{./graphs/classic_mapping.pdf}} \caption{
(a) The memory layout of a particular RAIMD-5 memory system for balanced
row-buffer locality and bank-level parallelism. (b)  The block layout of RAID-5
for comparison.  } \label{fig:raimd5_address_mapping} \end{figure*}


%\begin{figure}[!t] \centering
%\includegraphics[width=0.48\textwidth]{./graphs/stripe_parity_group.pdf}
%\caption{Relationship of stripes and parity groups in the linear address space.
%Seven data blocks and one parity block belong to the same parity group.  }
%\label{fig:stripe_parity_group} \end{figure}

Fig.~\ref{fig:raimd5_address_mapping} shows the detail of address mapping and
memory layout for a particular RAIMD-5 memory system of 1:7 parity-to-data ratio.  It
uses a stripe-based address mapping scheme to balance row-buffer locality and
bank-level parallelism.  A stripe is consecutive $K$ blocks in the physical
address space, where $K = 4$ by default.
Fig.~\ref{fig:raimd5_address_mapping_stripe} shows the layout of data stripes
and parity stripes, where $S$ represents a data stripe and \textit{PS}
represents a parity stripe.  For example, $\textit{PS}_0$ is the parity stripe for $S_0$ to
$S_7$, $\textit{PS}_1$ is for $S_7$ to $S_13$, and so on.  As the figure shows,
the data and parity stripes in any parity group are evenly distributed to all
devices, and thus there is no bottleneck for parity accesses.  The RAIMD-5
layout is revised from the conventional RAID-5 layout, as shown in
Fig.~\ref{fig:raimd5_mapping_conventional}, by moving the parity stripes down
to the bottom. This movement simplifies the physical-to-device address mapping.
It does not incur any performance or energy penalty in
a memory system; that may not be true to a storage system of mechanic disks.

With the above address mapping and memory layout, RAIMD can continue to use
conventional physical-to-device address mapping.  It incurs no extra delay on
data block access.  Parity block access, on the other hand, has an extra delay
because of the Parity Address Mapping (PAM) logic (Fig.~\ref{fig:mapping}), which
maps a data block address to the associated parity block address.  The logic may
involve shift, module, division, and multiplication logics and may take multiple
clock cycles to complete.  However, this extra delay is not on the critical path
of the program execution, and thus it has negligible impact on performance.
Thus, this framework allows a designer to use a complex address mapping to
achieve a desirable memory layout, e.g., for a different parity-to-data ratio, stripe
size, or having double parity blocks in a parity group.  A memory controller
may have multiple PAM logics to support different memory layouts.

\paragraph{Parity update mechanism.}

\begin{figure} \begin{subfigure}{0.45\linewidth} \centering
\includegraphics[height=3cm]{./graphs/raim_cache.png} \subcaption{Overview of
the parity update process.} \label{fig:raimcache} \end{subfigure} \hspace{1cm}
\begin{subfigure}{0.45\linewidth} \centering
\includegraphics[height=4cm]{./graphs/io_gating.pdf} \subcaption{Modification to
memory I/O circuit (circled).} \label{fig:io_gating_parity_update}
\end{subfigure} \caption{Parity update mechanism in RAIMD.}
\label{fig:parity_update} \end{figure}

The second critical issue in RAIMD is the increase of memory traffic for parity
block access.  RAIMD uses a novel parity update mechanism to reduce this
overhead, as shown in Fig.~\ref{fig:parity_update}.  We assume that the LLC is a
write-back cache, which is true for contemporary high-performance processors.
The LLC and the memory controller are extended to support a new request called
parity update, which is treated as a special type of write request.  Upon a
dirty write-back to the LLC, the XORed difference $\Delta D$ of the old data and
the new data is generated (steps \textcircled{1} and \textcircled{2}), and the
physical address of the associated parity block is generated
(Fig.~\ref{fig:mapping}).  A parity update request is then generated and
inserted into the write queue of the LLC.  The data block will be written back
to cache (step \textcircled{3}).  Independently, $\Delta D$ will merged into a
cache block for the parity block address if it exists; or written to a newly allocated
cache block if otherwise (step \textcircled{4}).  The parity block will cause a
parity block update to memory when it is evicted from LLC.

To serve a parity update request to the memory, the memory controller may issue
up to three memory commands, namely precharge, activate, and parity update.
Memory devices need to support a new command called parity update.  Its timing
is the same as the write command.
Fig.~\ref{fig:io_gating_parity_update} shows the circuit modification inside a
DDR5 x8 device to support the new command.  We have done a thorough circuit
analysis and concluded that the change does not affect the timing parameters
related to the write command.  That is because in modern memory access
scheduling, there is a sufficient time interval between the arrival of a write
command and the arrival of the write data such that the added circuit does not
increase the critical path delay (except for the delay of the XOR gate, which is
negligible).

With this parity update mechanism, parity read traffic is eliminated (except
during error correction), and parity write traffic is significantly reduced
because of caching.  Our performance evaluation confirms that parity traffic is
very moderate; the average increase of memory traffic is less than 5\% over the
baseline, and the average increase of cache misses is less than 2\%.  It is
worth mentioning that a similar idea called partial parity caching has been
proposed and studied to improve the performance and lifetime of flash-based
RAID-5 systems~\cite{flash_raid, partial_parity_cache}, but the implementation in
RAIMD is very different.

\paragraph{Evaluation Results}

We have evaluated the performance and energy efficiency of multiple RAIMD
variants using gem5~\cite{binkert2011gem5} and SPEC CPU2006 and CPU2017
workloads.  We focused on a particular RAIMD-5 configuration of 1:7
parity-to-data ratio based on x8 DDR5-6400 devices.  It uses a 32-bit error
detection code per 512-bit data block to detect errors.  The code is generated
and checked by the memory controller and stored in memory device, and therefore
the detection covers memory cells, bus, and I/O circuits.  The result shows that
RAIMD-5, which supports ChipKill protection, reduces memory energy consumption
by 25.9\% on average over a DDR5 baseline with no ChipKill protection.  It
incurs a moderate performance loss of 3.8\% on average, which is mostly caused
by increased data transfer time in memory access latency.  When compared to
unity ECC~\cite{unity_ecc}, a recent DDR5-based ChipKill design for servers,
RAIMD-5 reduces memory energy consumption by 69.4\% on average.  It is not a
surprise, because unity ECC has to distribute the content of any ECC word to ten
devices, and thus every memory access involves ten devices.  In general, parity
block based error correction is more efficient than ECC-based one for memory,
because the former can use only one device to serve a memory request.  The total
storage overhead ratio of the above RAIMD-5 configuration is 20.54\%, with
14.29\% for parity and 6.25\% for error checking code.  By comparison, the
storage overhead ratio of unity ECC is 31.25\%, with 25\% for extra ECC devices
and 6.25\% for on-die ECC.

A RAIMD system can be dynamically configured during system booting time as long
as the PAM logic supports the needed address mapping. Our evaluation shows that
a RAIMD-6 configuration, which supports double ChipKill protection, reduces
memory energy consumption by 23.7\% and incurs a performance loss of 4.3\% when
compared to the DDR5 baseline.  A server computer may opt to use RAIMD-6 for the
extra reliability.  A consumer-level computer, on the other hand, may use
RAIMD-5 with a smaller parity-to-data ratio to lower the storage overhead and
thus reduce the cost.

\subsubsection{Decoupled Memory Access Scheduling}

The second preliminary study focused on decoupled memory access scheduling
(decoupled scheduling in short thereafter) to support a heterogeneous memory
system.  In a conventional memory system, the memory controller conducts
centralized scheduling of memory operations, including precharge, activate,
read/write, auto-refresh, and other operations.  All memory devices behave as
slave devices to the memory controller.  This was necessary because it avoids
memory bus conflict and simplifies device circuit design.  However, it is severe
limitation to the adoption of emerging, non-volatile memory technologies,
including PCM (Phase-Change Memory), STT-RAM (Spin-Torque-Transfer RAM), ReRAM
(Resistive RAM), and MRAM (Magnetic RAM).  They are promising
alternatives to DRAM, but they behave differently and require careful management
for longevity, reliability, performance, and/or energy efficiency.  It is almost
infeasible for a memory controller, which is an integrated part of a processor,
to support those diverse technologies using centralized scheduling of memory
operations.

We have proposed a form of decoupled scheduling called POMI (POlling-based
Memory Interface)~\cite{Trung:pomi}.  It is an improved design over UniMA, an earlier
design of decoupled scheduling by the co-PI~\cite{Fang:pact11}.  It is also closely related
to Twin-Load~\cite{Cui:memsys16}, which supports a hybrid memory system of two types of
memory modules with different latencies.  POMI utilizes a device-specific buffer
chip per memory module to implement a polling based memory bus protocol.  The
buffer chip receives read and write requests from the memory controller and
performs local scheduling of device-specific operations, including precharge,
activate, read/write. The memory controller conducts global scheduling of read
and write requests and no longer tracks the status and ongoing operations at
every memory module.  The design not only reduces the complexity of the memory
controller, but more importantly allows it to interoperate with memory modules
using diverse technologies.

Decoupled scheduling, however, requires a new mechanism to avoid bus conflict. A
new polling-based is proposed to communication between memory modules and the
memory controller.  When receiving a read request from the memory controller,
the buffer chip of the target memory module sends a feedback with an estimated
ready time.  The memory controller schedules the bus usage, and it may
send a transfer signal, which indicates bus availability, to a memory module
based on its estimated completion time.  The latter may perform data transfer
after receiving the transfer signal, if the data is ready.  Otherwise, if the
approximation is wrong, it sends another feedback with a new estimated
completion time.  The evaluation shows that POMI can actually improve the
average performance of memory-intensive workloads by 3.7\% for a hybrid memory
system of both DRAM and PCM.  We found that polling adds an average time of 4.1~ns
to memory latency, but it is more than offset by increased parallelism in
memory modules because of local scheduling.

The POMI design shows the potential of decoupled scheduling, but it may not be
migrated to RAIMD directly.  In RAIMD, local scheduling will be implemented into
memory devices, not in a separate buffer chip.  Memory device design is very
sensitive to cost, and thus a more lightweight design is needed.

\subsection{Decoupled Scheduling in RAIMD}

\paragraph{The Mechanism.}

\begin{figure}[!t] \centering
\includegraphics[width=0.7\textwidth]{./graphs_new/device.pdf}
\caption{High-level view of the internal structure of a memory device with local scheduling.}
\label{fig:device}
\end{figure}

Decoupled scheduling is critical to the proposed RAIMD architecture.  It allows
the local enhancement for error protection and support for heterogeneous memory
technologies.  Fig.~\ref{fig:device} shows a high-level view of a generic
memory device in RAIMD with local scheduling and enhancement circuits.  The
memory controller issues three types of commands to a memory devices, namely
activate, read/write, and inquiry. The former two commands play the same roles
as in DDR4/DDR5 memories: an activate command moves a row of data to the row
buffer, and a read/write command causes the read/write of a memory block from/to
the row buffer.  The read/write command contains a bit to indicate if the memory
bank should be closed afterwards.  The inquiry command is new, with which the
memory controller may inquire if an activation has completed, if read data is
ready in the bank, if the bank is ready for incoming write data, and other
conditions.

Local scheduling is lightweight: its essential role is to decide the timing of
memory operation for each command at each bank.  It also responds to the inquiry
command.  The schedule logic maintains the status of each bank, enforce memory
timing requirements (which decides the start time of an operation), and tracks
the completion of the ongoing operation.  In other words, the memory controller
schedules the commands to devices, and a device schedules the timing of the
corresponding operations.  The memory controller schedules the bus usage to
avoid bus conflict.  It may only send a read command to a device if the data is
ready in the row buffer, and a write command if the row buffer is ready to take
on the incoming data.  Finally, local scheduling logic also schedules other
operations not directly related to memory requests, including refresh for DRAM
devices.

The lightweight design ensures the cost effectiveness and operational efficiency
of memory devices.  There is no additional data buffer or queue, which would be
expensive and hurt energy efficiency.  Note that the design is different from
POMI: the buffer chip of POMI is a separate chip and schedules memory commands
to serve requests. In RAIMD, the latter is still the role of the memory
controller.  The local scheduling logic will be implemented as a small state
machine. The cost is small in today's memory devices, whose I/O circuits are
already a significant part of the chip area.  The memory controller is also
simplified, as it no long tracks the status and ongoing operations of all memory
devices.  That can be costly if a memory controller has to support various
memory technologies.  It is worth mentioning that memory controllers are
commonly embedded into processor chips, and thus this issue actually affects the
interoperability of processor chips and memory devices.

Decoupled scheduling may have a moderate impact on performance because the
inquiry and feedback process will add a delay to memory access latency.  Our
preliminary study on decoupled scheduling, however, indicate the impact is can
be offset by increased memory parallelism, and decoupled scheduling may even
improve performance.  With decoupled scheduling, a memory command can be sent to
the target memory device as soon as it is decided and the command/address bus is
ready.  With rigid scheduling, the memory controller must make sure the device
will be ready for the operation when it receives the command (there are a number
of timing constraints), otherwise it has to delay the command.  The proposed
research will give a quantitative answer to this question.

\paragraph{Local enhancement circuits.}  With decoupled scheduling, memory
devices can have device-specific and technology-specific enhancement circuits
for various purposes.  For DRAM devices, those may include efficient DRAM
refresh, Rowhammer defense, and local error correction.  For NVM devices, those
may include wearout management, cell fault tolerance, and enhanced error
correction.  Those issues are not directly related to reliability, but they are
important and can be addressed in the RAIMD architecture.

DRAM refresh is increasingly a concern to the performance and energy efficiency
of high-capacity DRAM devices because more and more DRAM rows have to be
refreshed in DRAM refresh cycle, which is 64 ms under normal temperature by the
JEDEC standard.  A study has shown that, for memory-intensive workloads, DRAM
refresh energy consumption can reach 50\% of the entire memory energy
consumption, and the system performance may degrade by up to
30\%~\cite{Bhati:isca15}.  An promising approach is retention-time aware DRAM
refreshing~\cite{hamamoto1998retention,liu2012raidr,Liu2013,Patel2017}, which refreshes strong and weak rows at different rate.
Most DRAM rows are strong rows, with retention time much longer than 64 ms.  The
problem in a conventional memory system, however, is that the memory controller
has to manage the extra complexity and schedule fine-grain refresh commands to
all memory devices.  In RAIMD, it can be efficiently implemented as a local
enhancement circuit, which can maintain a record of a small number of week rows
(by profiling at booting time) and issue the fine-grain refresh commands
locally.  Because of decoupled scheduling, those locally issued refresh commands
will not cause any hard conflict with the incoming commands from the memory
controller.

DRAM Rowhammer attack~\cite{kim2014flipping,kim2020revisiting} is increasingly a security concern
as DRAM manufacturing moves to more advanced fabrication technology.  In this
attack, an adversary may flip bits in a DRAM row by frequently accessing its
neighbor rows to cause disturbance errors.  Many solutions have been
proposed~(\cite{aweke2016anvil,you2019mrloc,wang:iccd21} for a few examples), but they incur non-trivial performance
and energy overhead and are still not strong enough.  A new extension of JEDEC
DDR5 standard has introduced an optional Rowhammer mitigation framework called
Per Row Activation Counting (PRAC)~\cite{JEDEC:2024}, which embeds per-row
activation count in DRAM device and uses a special RFM (refresh management)
command.  However, a recent study has found the performance loss of PRAC can
reach over 80\% for DRAM devices of very low disturbance
threshold~\cite{Canpolat:hpca25}.  The cause is mostly the inefficiency of the
memory controller issued RFM command: it may block a memory bank for hundreds of
nanoseconds.  In RAIMD, the refresh management can be done by a local
enhancement circuit with much higher efficiency.  A local enhancement circuit
can issue the RFM command, and each instance of refresh management can be
fine-grain and take much shorter time.  Again, because of local scheduling,
locally issued RFM commands will not cause any hard conflict with memory
controller issued commands.

Wearout management is an important issue to NVM devices.  Many such designs have
been proposed.  For example, an NVMD may contain an SRAM cache to reduce NVM
writes, bit rotating design for wearout leveling, fault tolerant circuits to
handle hardware errors, among many other design
ideas~\cite{Lee:ISCA09,Zhou:ISCA09,Qureshi:ISCA09,Qureshi:HPCA10}.  In a
conventional memory system with rigid scheduling, it is infeasible to have a
memory controller to support various NVMDs of those techniques, which requires
the memory controller understand the internal of each device type.  In RAIMD,
the wearout management can be implemented as a local enhancement circuit, which
schedules those operations when needed.  Those operations, when ongoing, may
delay activate and read/write operations, but again they will not cause a hard
conflict in the scheduling.

We will study the integration of those designs into RAIMD as local enhancement
circuits.  Our focus is not on the circuit designs themselves, but on the
efficiency of the integration and the performance impact.  Some of those designs
may cause a high variation to operation completion time, which is a concern to
the overall performance of the baseline RAIMD.  We continue to propose neural
model based prediction to address this issue.

% Write a paragraph about DRAM refreshing and Rowhammer attack

% Write a paragraph about rowhammer attack

% Write a paragraph for NMV fault tolerance

% TODO discuss more about rowhammer, PRAC, and RFM

\subsection{Neural Model Based Prediction for Decoupled Scheduling}

The inquiry and feedback process of decoupled scheduling adds a delay to memory
access latency, which we predict is a concern for memory devices with advanced
enhancement circuits.  We call them advanced devices thereafter.  It is unlikely
a concern for simple devices as our preliminary study~\cite{} shows.  However,
advanced devices may have much higher timing variation in their operations than
simple devices.  For example, a DRAM device with Rowhammer defense will trigger
internal target refreshes, which may block normal operations (activate,
read/write, and associated precharge).  An NVM device with wearout leveling
circuits may trigger internal data movement, which may block normal operations
temporarily.  Even for simple devices, internal DRAM refresh may cause delay to
normal operations, which was not considered in the preliminary study as it uses
centralized scheduling for DRAM refreshes.  As the result, the overhead from the
inquiry-and-feedback process will increase.  The impact will be quantified in
the proposed research, but new methods should also be investigated.

\begin{figure}[!t] \centering
\includegraphics[width=0.7\textwidth]{./graphs_new/neural_model.pdf}
\caption{A high-level scheme of NM-RTP (neural model based ready-time
prediction) for RAIMD.}
\label{fig:neural_model}
\end{figure}

To improve the accuracy of ready-time prediction (RTP), we propose to use neural
model based read-time prediction (NM-RTP) in RAIMD.  Neural model with online
learning has been used successfully in branch prediction~\cite{} and memory
prefetching~\cite{}.  We are also studying the use of a perceptron-based neural
model to improve the accuracy of memory prefetching using memory-side
information~\cite{}.  We believe a good neural model, particularly
perceptron-based, is promising to improve RTP accuracy, but it has to be done
carefully.  Fig.~\ref{fig:neural_model} illustrate the NN-RPT scheme that we
propose.  The ``Inquiry*'' signal includes activate, read/write, and inquiry
commands, for all of which the device will send back a feedback.  The feedback
includes an estimated ready time and a feature value, the latter of which will
be used as a device input to the neural circuit.  A memory device contains a
circuit for device-specific feature generation, which is expected to have a good
correlation with the actual ready time.  It should indicate how reliable the
estimated ready time is, and provide a hint about the actual ready time if the
estimate is wrong.  The neural circuit will learn from all the feedbacks from
the device, and then make a best prediction of the ready time, which assits the
memory controller to schedule memory commands.

The majority of the cost lies in the weight tables of the neural circuit.  It is
important that the neural circuit is located in the memory controller, not in
the memory devices.  On the other hand, the circuit of device feature generation
must reside in the device, and it should be customized for each device type so
that the feature value is highly correlated to the ready time.  The feature
generation is important.  For DRAM devices, it can be a hint for incoming or
ongoing DRAM refresh, Rowhammer defense operations, low-power mode and mode
switch, and others.  For NVM devices, it can be a hint for incoming or ongoing
wearout leveling operations, fault tolerance related operations, chance for buffer
hits (if a performance-enhancing buffer exits), and so on.  The whole scheme is
a form of distributed learning circuit design: the neural circuit in the memory
controller carries out most of the learning, and the circuit in the device
provides its best assistantship.  Finally, the neural circuit may also include
other inputs, either from the inside of the memory controller or from the
processor cores, as secondary features.  In the proposed research, we will
investigate the full detail the NM-RTP scheme, evaluate its performance, and
optimize it to improve the accuracy.

%\subsection{Holistic Error Protection and Reliability Measurement Methodology}
\subsection{Holistic Protection and Reliability Measurement Methodology}
\label{sec:holistic}

RAIMD provides a form of global protection: when an memory error is detected by
the error checking code, the content of the faulty block is reconstructed by the
blocks in other memory devices in the same parity block.  It can correct any
single-block error.  However, given the worsening raw memory reliable, the
probability of double-block error cannot be ignored, particularly for servers
that require very high reliability.  RAIMD-6 can correct double-block error, but
it increases the cost and is less energy-efficient than RAIMD-5.  After all,
triple-block error may also become a concern for servers, which cannot be
corrected by RAIMD-6.

RAIMD can be enhanced with local protection by using embedded ECC.  The PI has
studied embedded ECC for DRAM memory, in which ECC bits are stored with the
associated data bits in the same DRAM row~\cite{Chen:taco13}.  It is similar to
on-die ECC of DDR5 but is more advanced.  Embedded ECC has a unique advantage
over SECDED ECC memory: it can use very long ECC word size.  This is important,
because the storage efficiency of ECC will increase drastically.  Assume the
data part of the ECC word is 512 bits, and the BCH code~\cite{} is used as the
codec.  It takes only 10 ECC bits to correct a single-error, 20 ECC bits to
correct a double-error, and 30 ECC bits to correct a triple-error.  The storage
overhead ratio is 1.95\%, 3.90\%, and 5.86\%, respectively.  It is much lower
than the 12.5\% overhead of SECDED ECC using an ECC word of 64 data bits and 8
ECC bits, which can only correct single-error.  The BCH circuit can be part of
the local enhancement circuits.  It is worth mentioning that the ECC word size
is limited by the row buffer size, not the LLC block size, and thus the data
part of the ECC word can be increased to 1024 bits or longer for even better
storage efficiency, though the cost of the ECC circuit will also increase.
There is a large design space on local error correction, including the choice of
the codec, the ECC word size, the design of the ECC circuit, and the performance
and energy efficiency evaluation.  We will fully investigate the design space.

Furthermore, the global protection and local protection together form a type of
holistic error protection, and they can be balanced to maximize the protection
strength for a given storage budget.  An uncorrectable error will happen to
RAIMD-5 if and only if two or more blocks in the same parity group have memory
faults that cannot be corrected by embedded ECC.  Therefore, the holistic
protection is much stronger than a baseline RAIMD-5, and it can be more
efficient than RAIMD-6 in storage overhead and energy consumption.  If extremely
high protection is required, e.g.~for devices of weak raw reliability, RAIMD-6
with embedded ECC can also be used.  With the objective to maximize reliability
for a given storage budget, there is another large design space to explore,
including the choice of RAIMD-5 and RAIMD-6, the choice of parity-to-data ratio, and the
choice and configuration of embedded ECC.  We will also fully investigate those
issues.

Finally, we will study new reliability measurement methodology to evaluate and
quantify the reliability of RAIMD systems.  It is necessary for searching the
optimal configuration of the holistic protection, but it is challenging given
the complexity of RAIMD systems and the complexity of memory fault patterns.  It
is impractical to use either pure analytical method or pure Monte Carlo
simulation.  A recent study introduces a new methodology to conduct empirical
analysis of DRAM fault model; and has developed a fault and error simulator to
evaluate the effectiveness of existing ECC schemes~\cite{jun:micro23}.  The
simulation framework runs Monte Carlo trials to evaluate the reliability of ECC
schemes, using error patterns produced from the fault model.  It is unique in
that it actually executes the ECC function in the Monte Carlo trials, which
helps improve the evaluation accuracy for complex ECC schemes.  We plan to
extend this methodology and the simulation framework to evaluate and quantify
the reliability of RAIMD systems.  The challenge is that RAIMD systems are much
more complex than ECC schemes, particularly when local error protection is used.
Thus, the error patterns must be generated with a sound methodology; otherwise,
either the Monte Carlo simulation takes excessive time to finish, or the
reliability measurement will not be accurate.  We will study this methodology in
depth.

% The case here is different from RAID-5: In RAID-5, a faulty block loses
% all the data, but in RAIMD that's usually not the case. Therefore,
% the codec design space is different.

%TODO Check a detail list of memory errors.

\section{Related Work}

\paragraph{ECC and ChipKill Memories} As memory systems become more prone to
multi-bit errors, various techniques have been proposed and applied to protect
against memory errors.  Redundant chips can be added to the DIMM to provide
chip-level or rank-level
protection~\cite{IBM:whitepaper97,DELL:chipkill00,Intel:note02, ChipKill,
ddr5_fault_bound, unity_ecc}.  ChipKill has been proposed to correct single-chip
errors and detect double-chip errors~\cite{ChipKill} by distributing ECC bits
across multiple memory chips ensuring that a failure in a single chip affects
only one bit of any ECC word.  However, current schemes for chip-level
protection either come with high performance overhead or low energy efficiency.
A recent study proposes Unity ECC~\cite{unity_ecc}, which utilizes Rank-Level
ECC (RL-ECC) for DDR5 memory.  Two redundant x4 devices holding ECC bits along
with eight x4 devices holding data bits are connected to a sub-channel.  The
eight ECC bits along with the 32 data bits are transferred on the 40-bit wide
data bus to provide rank-level protection.  This will result in high storage
overhead, 32.8\% with x4 devices assuming on-die ECC bits take up 6.25\% as
redundant cells~\cite{unity_ecc}.  It also incurs at least 25\% energy overhead
since each DRAM access will access eight data chips and two redundant chips.
DDR5 has not yet received a widely accepted ChipKill-level ECC
proposition~\cite{Row_Hammer_Resilient_DRAM} due to the split channels.  We use
this design for comparison in our evaluation of RAIMD because it is a recently
proposed, DDR5-based design.

There are also non-conventional ChipKill correct solutions.  For example,
LOT-ECC~\cite{Udipi:ISCA12} provides ChipKill correct capability using
nine-device rank (assuming x8 devices), but with an increased ratio of ECC bits.
Virtualized ECC~\cite{Doe:ASPLOS10} supports flexible error-correcting codes and
stores data and ECC bits in separate locations.  As another ChipKill
alternative, Nair \etal propose XED~\cite{Nair:isca2016} that utilizes on-die
ECC to detect error and use DIMM-level parity to correct the error in a way
similar to RAID-3.

\paragraph{RAID-Like memory systems}

There have been several existing studies using redundant memory with RAID-like
layouts, such as Intel Memory Mirroring~\cite{intel:mirror} with RAID-1 layout,
IBM RAIM~\cite{IBM:12} with RAID-3 layout, and RAIM5~\cite{RAIM5:17} with RAID-5
layout. Intel Memory Mirroring stores a backup copy for the entire memory
channel~\cite{intel:mirror}.  IBM RAIM provides comprehensive error protection
against channel failures as well as bus failures and full DIMM and chip
failures~\cite{IBM:12}. However, such strong protection comes with very high
cost. It uses four memory channels with ECC and one channel with parity data,
which results in nearly 40\% storage overhead. Each memory read or write
accesses 256 bytes across five memory channels simultaneously, which has a
negative performance impact and significant energy overhead.  Furthermore, it
requires an advanced memory buffer (AMB) to handle the upstream and downstream
data. Another study named RAIM5~\cite{RAIM5:17} proposed an optimized RAIM
configuration, which incurs 17\% off-chip traffic overhead with 28\% DRAM energy
overhead.  Another work named XED~\cite{XED} uses RAID-3 layout.  It can recover
single-chip failure using the parity data stored in a ninth chip. The chip
failure is detected by the on-die ECC; and a pre-defined ``catch-word'' is sent
to the memory controller. As with the other designs, XED consumes more energy
than conventional memory.

\paragraph{Sub-rank organizations}

The baseline RAIMD design is closely related to sub-rank memory organizations,
such as Mini-Rank~\cite{minirank} and Multicore DIMM~\cite{SC09}, to achieve
high energy efficiency.  For instance, Mini-Rank~\cite{minirank} divides each
64-bit rank into multiple sub-ranks and activates only a sub-rank of devices for
each request.  The Mini-Rank design does come with a performance penalty because
the data transfer time for a memory block, typically 64 bytes, is reciprocal of
the sub-rank size.  Thus, the original Mini-Rank study (based on DDR3-1066)
found that dividing a rank into two mini-ranks, e.g., four x8 devices per
mini-rank, achieves the best energy-performance trade-off for most applications.
However, with the current trend of ever-increasing DRAM data rate, it is now
feasible to use a sub-rank of a single device, e.g., one x8 device per
mini-rank.  This change not only further improves energy efficiency but also
enables error protection by the RAIMD design.  Even though in Mini-Rank, ECC
bits can be used and are distributed across different mini-ranks as in RAID-5,
the ECC bits are not embedded with the data bits as on-die ECC, thus each memory
access will need to access extra chips in one mini-rank, which is not efficient
and wastes the memory bandwidth.  In RAIMD, we address this issue by using
parity data to form a RAID-5 like architecture, and on-die ECC to detect errors,
forming strong error protection and detection against chip failures, but also
achieving energy efficiency at the same time.

\paragraph{Earlier studies on decoupled memory access scheduling}

Regarding the memory access protocol design, a recent work called
Twin-load~\cite{Cui:memsys16} allows one processor to support both
direct-attached and buffer-on-board memory with minimal changes to the existing
\DDRx protocol by splitting one read request to two loads.  It changes the
access protocol slightly, but not as comprehensive as the proposed design.  In
an earlier study, the PIs have proposed a framework, called
UniMA~\cite{Fang:pact11}, to support heterogeneous memory modules in one system by extending the
function of bridge chip and offloading the scheduling of device-level operations
to each memory module.  It uses a token-based protocol for the bus communication
between memory controller and modules.  However, the bridge chip does not have
any advanced functionaries, and the access protocol is different and the
interface is more complex.


\section{Results from Prior NSF Support of Last Five Years}

The PI and co-PI have not been supported by NSF for the last five years.

%The PI has been mainly supported by NSF grants ``SHF: Medium: Collaborative
%Research: Architectural and System Support for Building Versatile Memory
%Systems", and ``SHF: Small: Enhancing Memory System Dependability by Integrity
%checking". The objectives are to understand the complex interaction between
%multicore processors and memory subsystems; to develop architectural and system
%support for versatile memory systems; and to improve memory reliability. Related
%publications can be found in~\cite{Mittal:vlsid13,Mittal:iccd13,
%Chen:iccd13,Lin:tc13,Fang:tc14,Chen:taco13,Mittal:tvlsi13,Chen:isca14,wang:cal19}.
%Partially supported by the grants, four PhD students have graduated
%and another three PhD students
%are working on their dissertations.

%The co-PI has been mainly supported by NSF grants ``CAREER: Scalable and universal
%architecture for next-generation memory systems" and ``SHF: Medium:
%Collaborative Research: Architectural and System Support for Building Versatile
%Memory Systems". The two projects are to develop framework for next-generation
%memory systems and architectural and system support for versatile memory
%systems.  The objectives are to improve memory performance and energy
%efficiency; and to develop the framework to support emerging, diverse memory
%devices.  Related publications can be found in~\cite{Fang:ics13, Lin:tc13,
%Fang:tc14, Pourshirazi:ipdps16, Pourshirazi:memsys17, Pourshirazi:date18,
%Pourshirazi:TODAES,Rafique:icpp18,Rafique:memsys19}.  With the
%support of the grants, the PI has graduated two
%PhD students and one master student, and
%is supporting two PhD students.

\section{Project Management Plan}

The PIs plan to carry out the research with the following three-year plan:

\begin{itemize}

\item {\bf Year 1.\ }

We will focus on decoupled scheduling and NM-RTP.  On decoupled scheduling, we
plan to do Verilog modeling for the scheduling logic in memory device to verify
its functional correctness and to evaluate its circuit cost, extend our current
gem5-based simulation for RAIMD to model decoupled modeling, and run simulation
to evaluate its performance impact. On NM-RTP, we plan to develop the neural
model for RTP, do Verilog modeling for the neural circuit in memory controller
and the feature generation circuit in device, verify their functional
correctness, and evaluate their circuit cost and performance.  We will prepare a
paper on decoupled scheduling in RAIMD for conference publication.

\item {\bf Year 2.\ }

We will focus on NM-RTP evaluation, local error correction with embedded ECC,
and reliability measurement methodology.  On NM-RTP, we plan to extend our
simulation to model the neural model and its integration into decoupled
scheduling, and optimize the neural model based on the evaluation result.  On
the latter two parts, we plan to develop new versions of embedded ECC, develop
methods to measure their reliability, and study the integration of global error
protection and embedded ECC.  We will also develop Verilog modeling for embedded
ECC circuits to evaluate its circuit cost and performance.  On publication, we
plan to publish two or three conference papers based on the results.

\item {\bf Year 3.\ }

We will focus on holistic error protection and the continuous development of
reliability measurement methodology.  On holistic error protection, we will
focus on its integration with embedded ECC and the balance of the two to
maximize the reliability for a given budget of storage overhead. We will also
develop new reliability measurement methods to guide and evaluate the
integration.  On publication, we plan to publish another two or three conference
papers and prepare for journal publications.  We also plan to finalize and
publish our Verilog modeling and gem5-based simulation code on GitHub.

\end{itemize}

\section{Education Plan}


The educational impact of this project will be significant from the following
aspects at the public university:

\begin{itemize}

\item{\bf New curriculum development.\ } The proposed research will
contribute to the current development of new graduate topic classes
related to memory systems at the Department of Electrical and Computer
Engineering at University of Illinois Chicago (UIC).  The findings
and simulation tools developed from this research will be introduced
in the classrooms in a timely manner.

\item {\bf Undergraduate participation.\ } We will encourage highly motivated
undergraduate students to participate in the proposed research. The College of
Engineering at UIC has a summer intern program that encourages undergraduate
students to work at research labs. We have been actively participating in this
program.  The students will be selected from undergraduate classes in computer
organization and architecture, where we are major instructors for those classes.

\item {\bf Woman and under-representative minority student participation.\ }
University of Illinois Chicago is one of the nation's most diverse public
research universities and is a federally designated Minority Serving
Institution.  Being a woman faculty, the co-PI is especially committed to
increasing the number of women pursuing advanced degrees in computer science and
engineering. The PI has been working with a woman PhD student in preliminary
study of RAIMD, have advised multiple undergraduate underrepresented minority
students.  We will continue to expand outreach efforts to broaden the
participation of underrepresented minority students in research and higher
education.

\end{itemize}

\section{Broader Impact, Dissemination Activities and Technology
Transfer}

The proposed research will lead to a new memory architecture to build memory
systems that are highly reliable, energy-efficient, and cost-effective.  It
works for all general-purpose computers, including servers, personal computers,
and mobile devices.  It is also promising for emerging memory technologies, as
those technologies tend to be less reliable than current DRAM technology.  In
the long run, the research will have a major impact on the design of computer
systems.  The proposed educational component will enrich computer system
education related to memory system design for both undergraduate and graduate
students.

The results of this project will be widely disseminated to the research
community as well as practitioners through (i) technical publications in high
quality journals/conferences and presentations at relevant conferences; (ii)
open sourcing of the source code of our simulation tools; and (iii) tutorials on
these techniques at conferences and companies that have related business.

This proposed research is of great practical interest to the industry.  Our past
research and designs on memory systems have been adopted in processors and
chipsets~\cite{zhang:micro00}; and we expect effective technology transfer for
the proposed research.

\newpage \renewcommand{\thepage} {\arabic{page}} \setcounter{page}{1}

\let\oldthebibliography\thebibliography
\renewcommand{\thebibliography}[1]{%
  \oldthebibliography{#1}%
  \setlength{\itemsep}{1.5ex} % Increases space between items
  \setstretch{1.3} % Increases line spacing within an item
}

\begingroup
\setstretch{1.2} % Adjust factor as needed (1.0 = single, 1.5, 2.0)

\bibliographystyle{latex8}
\bibliography{ref1,ref2,bib/proposal,bib/hpca18-submit,bib/isca14,bib/memguard2,
bib/memguard2b,bib/medium2,bib/mybib,bib/taco,bib/nas15,bib/related,bib/canary,
bib/iccd21}

\endgroup
\end{document}

% FIXME Is it still required this way?

\section{Project Management Plan}
% FIXME Check current requirement

\section{Education Plan}
% FIXME Check current requirement








-----------------------------------------------------------------------------


Furthermore, emerging memory technologies require much stronger support for
reliability at architecture/ microarchitecture level.  significantly different
from  memory organizations designed for DRAMs.  Fig.~\ref{organization}
contrasts their differences.  In the new organization, each memory device (chip)
embeds a light-weight memory access scheduling logic to schedule memory
operations for itself.  The memory controller issues memory access requests,
instead of memory op devices.  Each device operates independently, yielding high
memory concurrency.  The interface between the memory controller and memory
devices is revised to allow variable access latency.  This is a critical
difference from conventional memory organizations (such as DDRx), because it
allows a memory device to incorporate circuits to improve performance, energy
efficiency, reliability, and endurance without violating any timing
requirements.

This new organization offers several advantages over the conventional E
organization.  Since a single device serves a memory request, instead of a rank
of multiple devices, energy efficiency is significantly improved.  Reliability
tends to be a challenge with narrow memory ranks -- this new organization
essentially uses very narrow rank of one device per rank.  However, as we will
show in our preliminary result, it can actually offer very high reliability with
low cost.  More importantly, it has excellent extensibility to support new
memory technology and device designs, which is critically important to NVMs.
The complexity of the internal structures of NVM devices can be hidden behind
the new controller/device interface, which allows novel designs for improving
NVM performance, energy efficiency, and write endurance.  Additionally, it will
also help address the major challenges in contemporary DRAM systems, including
efficient DRAM refreshing and Rowhammer protection, by allowing efficient
mechanisms inside DRAM memory devices.

It is not that the new organization lacks its own challenges; however, those
challenges can be addressed with technology advances and designs.  First of all,
the new organization requires high-speed (data rate) memory devices; otherwise,
memory latency will increases substantially.  At this time, it is no longer a
concern.  In the near future, the data rate of mainstream DDR5 memory devices
will reach 6400 MT/s.  Our preliminary result shows that with that data rate,
the performance of 8-bit rank size is very comparable to that of 64-bit or
32-bit rank size, and the energy efficiency is much better.

The second challenge is how to support memory error protection with narrow rank.
Conventional schemes use (72, 64) ECC words, which does not fit narrow memory
ranks.  This is a known issue, and one of the PIs has proposed a solution for
Mini-Rank~\cite{chen:taco13}.  We further propose RAIM-D (Redundant Array of
Independent Memory Devices) for SSMD memory system\footnote{It is significantly
different from previous studies on RAIM~\cite{xx} (see Section~\ref{sec:related}
for detail).}.  It is similar to the RAID organization for disks, with further
optimization to avoid the small write issue of RAID-5.  In RAIM-D, each memory
device detects error by itself, and the correction is done at system level.
This new organization not only support SECDED (single error correcting double
error detecting) but also SDDC (single device data correcting).  Our preliminary
result shows that a SSMD memory system with RAIM-D may incur only xx\% of
performance overhead but xx\% improvement of energy efficiency.  The result is
consistent with the trend shown by previous studies~\cite{} on Mini-Rank and
sub-ranking.  With the high rate of today's device, it is now feasible to use a
rank as narrow as 8-bit.

The third and the major challenge is to maintain memory channel efficiency while
allowing operation timing flexibility.  Such a flexibility is the key to utilize
diverse memory technologies, or even just to hide the increasing complexities of
DRAM devices~\cite{}.  In such recent designs (some done by the PIs), a memory
device may send a negative feedback upon a memory command if it is not yet ready
for the requested operation, with or without extra feedback about the ready
time.  Those studies have shown that the approach is effective for DRAM
memories, yielding the same level of memory channel efficiency as conventional
DRAM designs.  However, an advanced approach is needed for NVM and mixed
DRAM/NVM memory systems because their latency variant is much higher than DRAM.
Such an approach may also support advanced DRAM designs to address DRAM
reliability, DRAM refresh bottleneck and Rowhammer attack.

Additional information can be embedded within the feedback to help the memory
controller retry at a right time.  Recent studies~\cite{}, including two of the
PIs, have shown that high channel efficiency can be achieved in DRAM memory
system of this type.  However, NVM devices with write leveling circuits and
fault tolerance designs may expose much higher timing variations than DRAM.  To
address this issue, the PIs propose to use perceptron-based neutral networks for
the memory controller to predict when a memory device will be ready for a given
operation.

% operation timing flexibility

% Flex timing Is it still command? Self operated devices? The operation is self
% scheduled.  For a request, MC sends activate commands, wait for the data to be
% ready, then sends data requests.  Now there are two commands, activate and
% data transfer (ACT and DT).  The memory controller knows the bank structure.
% It activates a bank for a row, and then sends column address.  For a simple
% device, this is OK.  For a complex device, a memory block could be cached --
% is that OK? But some devices have narrow banks.  In other words, if a bank is
% activated with a row, that doesn't mean the data of the whole row is all
% available.  The reason for this design is efficiency -- there is no need for
% data buffer, and a device can be almost as cheap as current devices.  On the
% other hand, a device can also be complex for optimization.  Alternatively: MC
% sends the full address.  The devices does scheduling... that is complicated.
% So, what is exactly self-scheduling?  Activate delay is variant, and
% read/write delay is also variant.  The device accepts command, not request,
% but its operations have varying delay.  Self scheduled operations?
% Self-operating? Command has flexible timing, and there is no precharge.  and
% some commands are self-scheduled. This is complex.

% Recap: Devices receive activate and transfer (read/write) commands. That's all
% commands supported except configuration.  The device decides timing.
% Self-timed access? Partially self scheduled is OK, but the name is too long.
% Activate command: Make a bank ready to read/write a row.  Transfer command:
% transfer data to a block (column) in the row.  One problem is that, what
% happens if a bank/row is activated, but not used?  This is not really a
% problem -- if MC activates another row of the same bank, then the MC should
% know.  And if the device decides to close the bank, the device should
% re-activate the same row when a block is requested.  Another issue is that the
% device doesn't really know if the block is ready or not.  This is less an
% issue: We can require that any bank has a buffer to hold a single block.  So
% this issue is fixed.  So, it really about flexible timing -- self-scheduling
% would be a misleading term.

% major challenge is the scheduling of bus.  Typically non-deterministic
% scheduling is an issue, but with advance of NN, we may use ML to improve the
% efficiency.

% Flex-timing devices: It is more about giving device designers the option to
% optimize devices Flex-timing: Activate has flexible Can data be sent first,
% and then activate? A single data buffer for data?  Reverse data transfer and
% activate? Why not reverse? Send data first, and then activate and write.  But
% that will make the device more expensive, but why not?  For above question: A
% "minimum" device may not have extra buffer at all -- that's like the current
% devices. It may reject a TX request unless the data is already in row buffer.
% But if a device has some additional buffer, it can take a TX request
% immediately until the buffer is full.

The major challenge of the new organization is that, without rigid timing, the
memory controller has to predict the completion time of DRAM operations.

-- actually, the latter have a rigid interface such that the memory controller
must schedule memory operations with precise timing to ensure correctness.

Equally important, this new organization is much more extensible to NVM
(non-volatile memory) devices than the conventional one.  NVM devices usually
have internal structures, for example write buffers and write-leveling
mechanisms, to improve performance and write endurance, and thus are not a good
fit for conventional memory organization of a rigid interface.

the conventional, rigid memory organization, and this is very compatible with
new types of memory devices, particularly those non-volatile memory devices.

\paragraph{Overview} In this project, the PIs will investigate machine learning
throttling techniques for heterogeneous memory systems with a memory access
protocol of flexible timing.  The research extends from preliminary studies of
memory access protocols of flexible timing~\cite{}, which is the key to
utilizing heterogeneous memory technologies including DRAM and NVMs
(Non-Volatile Memories).

Unlike a conventional memory system, the memory controller of a VDL memory
system does not track the resources and status inside each memory module.  It
allows interoperability between processors and a wide range of memory modules,
including those made by DRAM and NVM technology. This feature is critically
important to the adoption of NVM technologies, which have a wide range of timing
parameters and may have complex internal structures for performance, write
endurance, fault tolerance, and energy efficiency.  The preliminary studies of
the PIs, as proof of concept, have shown that a VDL protocol performs very
comparably to \DDRx protocols for DRAM memories.  Further research is needed for
heterogeneous memories, particularly on those made by NVMs.

% Discuss why machine learning is needed
In the proposed research, the PIs will study machine learning techniques to
improve the efficiency of VDL memory systems for a wide range of memory modules,
made by DRAM, NVMs, and with complex internal structures.  {\bf In such systems,
it is critical to use ML techniques for the memory controller to gauge the
internal status of memory modules and then predict the data latencies of pending
requests}.  A conventional design, which precisely tracks each memory module's
internal status, will no longer be feasible because of the mounting complexity.
A VDL memory controller, as proposed, uses machine learning to gauge imprecisely
the internal status of each memory module and issues memory commands
accordingly.  It is the imprecise nature of the design that makes it possible
for a VDL memory controller to work with a wide range of memory modules, without
the need to know precisely their internal structure and status.  Consequently,
machine learning plays a critical role in the timeliness of memory commands and
thus the efficiency of the VDL protocol.

% Discuss some specifics of ML
The PIs further propose a unique design of distributed ML framework for VDL
memory systems.  In this framework, memory modules provides rich ML feedback per
memory command to facilitate machine learning inside the memory controller.
Those feedbacks give the memory controller {\em hints} about the internal status
of memory modules that are strongly correlated to memory efficiency.  The
feedback is memory module specific and thus can be customized by the module
designers.  Coupled with machine learning techniques, such a memory controller
can issue memory commands in an efficient manner without the precise knowledge
of memory module internal status.  With a high-quality ML model, a VDL memory
system may approach or even surpass the efficiency of a conventional memory
system of rigid timing, and meanwhile support heterogeneous memory modules.

% In my thought, the feedback includes retry or not (which is a major feedback),
% plus hints for efficiency. A question here: how about energy efficiency?  The
% hints can be more useful than the retry feedback: ML may not reliably improve
% retry efficiency, but may improve the other metrics.

% TODO Here I need a picture.

% TODO here I need a major revision of the following paragraph
The feedback includes a field that will be used as a feature for the ML circuit
of the memory controller.  [Here I must write something about the latency
feedback.]  With this input, the ML circuit predicts the correct timings of
memory command for each memory request.  Consequently, the memory controller do
not have to issue memory commands with strict timing -- it only has to make the
best prediction with the help of ML.  If a prediction is incorrect and the
memory command cannot be served, the memory controller will get a negative
feedback and retry the command at a later time.  This framework will drastically
reduce the complexity of the memory controller: instead of tracking the status
of all hardware structures of memory devices, which is infeasible for advanced
NVM devices, the memory controller only has to gauge the status using ML
modeling.

[In POMI design, the device may send back latency information. This can be added
to this protocol: Some binary combinations are reserved for certain feedback,
where the others are feature.]

[What feedback? First, tell if an operation can be done or not. Second,
Additional information about the feature.] [Scenario 1: Normal memory access,
sends precharge, activation, and then data transfer.  Device has latches for row
address, bank/bank group, and column address.  buffers may or may not be ready,
because device has some control of the buffers. MC sends row/bank/bg address
first, gets a feedback. Then feedback's meaning depends on device. For example,
it could be that buffer is ready, not ready; or it can have multiple means, if a
NVM device has additional cache/buffer.  The feedback may also say the row is
not ready and a retry is needed. Then, after some waiting, the MC sends column
address to get the data.  The feedback first tells if a retry is needed.
Additional feedback tells the feature.]


[Draw a diagram here]

[The memory controller still maintains address mapping.  That cannot be done
using machine learning. ML decides when to send memory command.]

[The feedback should also include information such as if a memory bank is ready
for the command and so on]

[Two things to say: Data latency is variable, and the memory controller doesn't
have precise knowledge of what's inside a memory module.]

In this research, the PIs will investigate ML techniques to improve the
efficiency of variable-latency memory access protocols for heterogeneous memory
systems, including DRAMs and various NVMs.

- ML is used to predict when data is ready - Additional information is used to
help ML prediction

In the preliminary studies, the PIs have proposed new memory access protocols
with carefully relaxed timing requirements, which is necessary for heterogeneous
memory systems.  [Define heterogeneous memories?]  In the proposed research, the

A key challenge in this approach is to address a potential loss of efficiency
from the flexibility of the memory access protocol.  In this research, the PI
proposes new ML throttling techniques to improve the efficiency. [Define the
goal of the research: Very high efficiency,  complexity]

The complexity of memory subsystem has been increasing substantially to meet the
ever-increasing demands for higher capacity, bandwidth, throughput, reliability,
and energy efficiency.  Furthermore, non-volatile memories are poised to
challenge DRAM, the dominant technology for the past decades.  This transition
will have far-reaching impacts on computer systems design [expand to OS,
security, and so on].  However, it also substantially increases the complexity
of memory subsystem design, which is under high pressure to meet the increasing
demands from data-intensive workloads and must balance between performance,
energy efficiency, and reliability.

[Need to revise this part ...]

A core issue to this problem is the rigid design of existing memory access
protocols, with \DDRx protocols being the dominant design.  A memory access
protocol articulates how a memory controller communicates to memory devices,
e.g., DDR4/DDR5 DRAM devices (chips).  The current \DDRx protocols evolve from
the SDRAM (synchronous DRAM) from 1990s [double check the date], with drastic
improvements of data rate.  However, like in SDRAM, they are still rigid
protocols, in which memory devices are slave devices that respond to memory
commands, e.g., precharge, row activation, and data read/write, in a very strict
timing requirements.  [Give an example of timing.]  The strict timing
requirement was necessary, in the era of DRAM memory, to minimize the cost of
DRAM devices and meanwhile maintain the correctness and efficiency of memory bus
transactions.

However, in the new era of heterogeneous memories, the strict timing requirement
has become a severe constraint.  It leaves no room for a memory device to adopt
any design improvements that affect the latency of memory operations.  For
example, [PCM write leveling, Rowhammer attack, and so on.] ....  Thus, many
creatives design ideas proposed so far may not be adopted because they will
violate the strict timing requirements.

A [add cap] flexible access protocol for main memory, however, must make very
careful trade-offs between flexibility and efficiency.  Otherwise, it may lead
to significant loss of performance or energy efficiency [I want to talk about
latency and bandwidth].  For example, ...

The PIs will investigate ML-throttled memory access scheduling for memory access
protocols with relaxed timing requirements.  The PIs have studies such protocols
in previous studies~\cite{}, for DRAM and NVM [double check on NVM], and found
that they can perform comparably or even better than conventional

The PIs have found in a preliminary study~\cite{} that ML is very promising in
reducing the inherent overhead of a flexible memory access protocol... [Write
about preliminary results]

The preliminary study was based on DRAM devices.  Research on a flexible memory
access protocol for heterogenous NVM systems is much more challenging, and much
more important, because of diversity of NVM devices. ... [Discuss timing
variants, write levering, error tolerance, ECC, but avoid redundance.]
